\chapter{Brugada Syndrome}

\section*{Definition and Overview}
Brugada syndrome is a rare, inherited heart rhythm disorder in which the electrical system of the heart is abnormal, increasing the risk of dangerous ventricular arrhythmias that can cause fainting, seizures, or sudden cardiac arrest. The heart is structurally normal in most affected people, and the problem lies in the electrical signals themselves. The condition often becomes apparent in adulthood, and many affected individuals remain unaware of it until an electrocardiogram (ECG) is performed or a family member is diagnosed. Diagnosis and long-term care require a specialist cardiology team, and management is aimed at reducing the risk of life-threatening arrhythmias.

\section*{Epidemiology}
Brugada syndrome is uncommon, affecting roughly 1 in 2,000 to 5,000 people in the general population. It is more common in men than in women, with symptomatic disease occurring around eight to ten times more often in males. The syndrome is found worldwide but is most frequently diagnosed in Southeast Asia, where it may account for some cases of sudden unexpected death in young men during sleep. Symptoms typically first appear in the 30s and 40s, although the characteristic ECG pattern can be present in children. The true prevalence may be higher than reported, because many carriers never develop symptoms.

\section*{Aetiology and Pathophysiology}
Brugada syndrome results from inherited faults in genes that control the flow of sodium and other ions across the membranes of heart muscle cells.

\begin{itemize}
    \item \textbf{Sodium channel dysfunction:} mutations, most commonly in the \textit{SCN5A} gene, impair the sodium current that drives the early phases of the cardiac action potential
    \item \textbf{Right ventricular vulnerability:} the electrical abnormality is most pronounced in the right ventricle, where cells recover at slightly different rates, creating a dispersion of repolarisation
    \item \textbf{Arrhythmia generation:} this electrical imbalance can trigger ventricular fibrillation, a chaotic rhythm in which the heart cannot pump blood effectively
    \item \textbf{Genetic and acquired causes:} a specific faulty gene is identified in roughly a quarter of families, but in many cases no genetic cause is found; some cases appear without any known family history
    \item \textbf{Triggers:} fever, certain medications, alcohol, and electrolyte disturbances can unmask the ECG pattern or provoke arrhythmias
\end{itemize}

\section*{Clinical Features}
Many people with Brugada syndrome remain entirely asymptomatic, while others experience symptoms that typically begin in adulthood.

\begin{itemize}
    \item Fainting (syncope), often occurring at night, during rest, or during sleep
    \item Palpitations or an awareness of the heart racing
    \item Sudden cardiac arrest, which is sometimes the first manifestation of the condition
    \item A characteristic ECG pattern that may appear and disappear over time
    \item Seizures in some people, which are actually caused by the underlying cardiac rhythm disturbance
    \item Symptom provocation by fever, certain medicines, alcohol, or large meals
\end{itemize}

\section*{Differential Diagnosis}
The ECG changes of Brugada syndrome can be mimicked by other conditions, and fainting has many possible causes.

\begin{itemize}
    \item \textbf{Acquired causes of the ECG pattern:} certain medications (including some antiarrhythmic and psychotropic drugs), electrolyte imbalance, myocardial ischaemia, and conditions compressing the chest
    \item \textbf{Vasovagal syncope:} simple fainting triggered by stress, standing, or heat
    \item \textbf{Other inherited arrhythmias:} long QT syndrome and catecholaminergic polymorphic ventricular tachycardia
    \item \textbf{Structural heart disease:} hypertrophic cardiomyopathy and other conditions that predispose to arrhythmia
    \item \textbf{Seizure disorders:} epilepsy misread as rhythm-related collapse
    \item \textbf{Fever and heatstroke:} severe febrile illness can produce transient ECG changes
\end{itemize}

\section*{When to Seek Medical Care}
If a family member has been diagnosed with Brugada syndrome, or if there is a personal or family history of unexplained fainting or sudden cardiac arrest, seek a cardiology assessment. Before starting any new medicine, ask a doctor or pharmacist whether it is safe in Brugada syndrome.

\textbf{Call 000 (or local emergency services) for:}
\begin{itemize}
    \item Fainting, especially during rest or sleep
    \item Chest pain, palpitations, or shortness of breath with collapse
    \item Any witnessed cardiac arrest -- start CPR and ask someone to call 000
    \item Sudden collapse with no clear cause
\end{itemize}

\section*{Diagnosis and Investigations}
The diagnosis is made by a cardiologist using a combination of tests.

\begin{itemize}
    \item \textbf{Resting ECG:} may show the characteristic coved ST-segment elevation in the right precordial leads
    \item \textbf{Provocative drug testing:} administration of a sodium channel blocker under continuous monitoring to unmask the pattern if it is not visible at rest
    \item \textbf{Echocardiogram:} to confirm that the heart structure is normal
    \item \textbf{Genetic testing:} a blood test to look for known pathogenic variants, which helps confirm the diagnosis and enables cascade testing of family members
    \item \textbf{Ambulatory ECG monitoring:} a Holter monitor worn for 24 hours or longer to capture rhythm during normal activity
    \item \textbf{Electrophysiology study:} occasionally used to assess the risk of dangerous arrhythmias in selected patients
\end{itemize}

\section*{Management}
Treatment aims to prevent dangerous ventricular arrhythmias and sudden cardiac arrest.

\begin{itemize}
    \item \textbf{Implantable cardioverter-defibrillator (ICD):} a device placed beneath the skin that detects and terminates dangerous rhythms; it is recommended for people who have survived a cardiac arrest or suffered arrhythmic syncope, and considered in other higher-risk individuals
    \item \textbf{Medication:} drugs such as quinidine may reduce the risk of arrhythmia in selected people who cannot have an ICD
    \item \textbf{Avoidance of triggers:} prompt treatment of fever, avoidance of alcohol, and avoidance of medicines known to worsen the condition
    \item \textbf{Lifestyle planning:} a written plan for managing fever and illness, and for what to do if symptoms occur
    \item \textbf{Family screening:} clinical and genetic testing of relatives to identify others at risk
    \item \textbf{Genetic counselling:} to explain inheritance patterns and reproductive implications
\end{itemize}

\section*{Complications}
The principal complications relate to the arrhythmia itself and its treatment.

\begin{itemize}
    \item Ventricular fibrillation and sudden cardiac arrest, the most serious risk
    \item Recurrent fainting, with the risk of injury from falls
    \item Inappropriate shocks from the ICD, which can be distressing
    \item Device-related complications, including infection, bleeding, or lead malfunction
    \item Anxiety and psychological distress for affected people and their families
\end{itemize}

\section*{Prognosis and Outcomes}
The outlook varies widely between individuals. People with no symptoms and a normal resting ECG generally have a low risk of dangerous rhythms, whereas those who have already fainted or survived a cardiac arrest are at substantially higher risk. With an ICD in place, the risk of death from the condition is greatly reduced, and many affected people live long, active, and normal lives. Because there is no cure, lifelong follow-up with a specialist heart team and careful attention to avoiding triggers are essential.

\section*{Prevention and Self-Care}
The condition itself cannot be prevented because it is inherited, but the associated risks can be actively managed.

\begin{itemize}
    \item Seek cardiology review if a relative is diagnosed, and consider genetic testing
    \item Treat fever promptly with paracetamol and cooling measures, since fever can provoke arrhythmia
    \item Avoid alcohol and medicines known to worsen the condition, checking with a pharmacist or doctor before starting any drug
    \item Carry a written summary of the diagnosis, medications, and emergency contacts
    \item Ensure family members and friends know how to perform CPR and use an automated external defibrillator
    \item Attend all cardiology appointments and keep the ICD device checked regularly
\end{itemize}

\section*{Sources and Further Reading}
The following organisations provide reliable information on Brugada syndrome for patients, families, and healthcare students.

\begin{itemize}
    \item NHS. \textit{Brugada syndrome: symptoms, diagnosis, and treatment.} https://www.nhs.uk/conditions/brugada-syndrome/
    \item Mayo Clinic. \textit{Brugada syndrome: overview and management.} https://www.mayoclinic.org/
    \item MedlinePlus. \textit{Brugada syndrome health information.} https://medlineplus.gov/brugadasyndrome.html
    \item MSD Manual. \textit{Brugada syndrome: professional overview.} https://www.msdmanuals.com/
    \item National Heart, Lung, and Blood Institute. \textit{Inherited arrhythmia syndromes and sudden cardiac death information.} https://www.nhlbi.nih.gov/
    \item SADS Foundation. \textit{Support and education for sudden arrhythmia death syndromes.} https://sads.org/
\end{itemize}
